\section{Completed generic training experiment}
\label{sec:t41-results}
The owner-authorized deadline experiment used six attempted configurations at
30 epochs, followed by a 100-epoch untuned control and a 100-epoch confirmation
of the best completed search configuration. All runs used seed 17 and the same
verified general YOLOv10-M initialization. No confirmation continued from a
search checkpoint. Specific-task training, repeated seeds and the original
twenty-attempt protocol were deferred explicitly.

Training ran on one NVIDIA H100 NVL through Slurm, with four CPUs, 32 GB host
memory and a four-hour limit per allocation. The isolated Python 3.12 environment
used the pinned THU-MIG YOLOv10 fork, FP32, batch 32 and two loader workers.
Ali's AP@[.50:.95] validation objective evaluated all 2,445 validation images
after each epoch. Only validation determined search and checkpoint selection.
Exact initializer and selected checkpoint SHA-256 values are recorded in
\file{docs/t41-results.md}.

\begin{table}[htbp]\centering\small
\caption{Actual generic study accounting. AP is validation AP on a 0--1 scale.}
\label{tab:t41-search}
\begin{tabular}{llll}
\toprule Attempt/stage & Slurm job & Outcome & Selected AP\\\midrule
Search 0 & 755 & Failed at epoch-6 validation & Undefined\\
Search 1 & 758 & 30 epochs complete & 0.868115\\
Search 2 & 759 & 30 epochs complete & 0.861131\\
Search 3 & 762 & 30 epochs complete & 0.862620\\
Search 4 & 764 & 30 epochs complete; selected & 0.869853\\
Search 5 & 767 & 30 epochs complete & 0.861840\\
Control confirmation & 770 & 100 epochs; selected epoch 68 & 0.877062\\
Candidate confirmation & 819 & 100 epochs; selected epoch 84 & 0.874833\\
\bottomrule
\end{tabular}
\end{table}

Attempt 0 failed because the fork clipped positive boxes wholly outside an image
to zero area. A scoped bridge retains the original positive geometry for those
predictions, without dropping detections or inventing extents. Diagnostic job 756
and qualification job 757 verified the fix and unchanged earlier predictions/AP.
The failed attempt remains counted with an undefined objective; its diagnostic
score is excluded from search selection. No interrupted optimizer continuation
is implemented or claimed.

Successful search and confirmation allocations totaled 9 h 57 min 06 s;
including the failed search gives 10 h 08 min 59 s, excluding pilots and
diagnostics. The confirmations took 2 h 47 min 35 s and 2 h 46 min 49 s, each
below four hours. The candidate's confirmation validation AP was 0.2228
percentage points below the control. This small search can remain within the
sampler's startup phase; it is not evidence of TPE convergence. Independent CPU
verification rehashed 350 successful epoch prediction records and 237 bound
runtime files, and exactly reproduced both selected validation scores.

\section{Frozen published-test comparison}
The two checkpoints, code, test annotations, all test image hashes and settings
were frozen before Slurm job 820. The run completed on 23 November at 08:45:31
UTC+4 in 1 min 27 s. Both checkpoints were evaluated on all 5,136 published
test images and 3,204 annotated boxes, with category 0 mapped to Explosive.
The test set was used only for this final comparison; it did not select a model,
checkpoint, threshold, seed or hyperparameter. No settings were changed afterward.

Inference used batch one, FP32, image size 640, score floor 0.001 and maximum
300 detections. Official pycocotools 2.0.11 computed COCO box metrics. A separately
frozen operating point used confidence 0.25 and IoU 0.50, without test-based
threshold tuning. The full predictions were saved in original-pixel COCO xywh
coordinates. The qualified geometry policy remains active; neither final model
produced an outside-box event. A separate CPU process reproduced all metrics
and operating-point counts exactly from the saved predictions.

\begin{table}[htbp]\centering\small
\caption{Frozen generic test results; AP/AR and precision/recall are percentages.
These are project measurements, not published-paper scores.}
\label{tab:t41-test}
\begin{tabular}{lrr}
\toprule Metric & Control & Search candidate\\\midrule
AP@[.50:.95] & 41.5886 & 48.0502\\
AP50 & 63.4039 & 66.5397\\
AP75 & 42.2905 & 52.3537\\
AR@100 & 49.2603 & 53.2709\\
AP small / medium / large & 33.49 / 35.87 / 47.47 & 27.05 / 43.57 / 53.57\\
Precision at confidence 0.25 & 89.0028 & 95.0830\\
Recall at confidence 0.25 & 48.7516 & 48.2834\\
True positives / false positives & 1,562 / 193 & 1,547 / 80\\
Missed annotated boxes & 1,642 & 1,657\\
False positives on annotation-empty images & 55 & 25\\
\bottomrule
\end{tabular}
\end{table}

The candidate has 6.4616 percentage points higher test AP, but lower small-object
AP and slightly lower fixed-threshold recall. Its validation confirmation was
also lower than the control. These are descriptive results for a preselected
pair, not a basis for choosing a deployment model after seeing test scores.
Both models miss approximately half the annotated boxes at the stated operating
point. An annotation-empty image is not assumed benign; no operational safety
or clearance claim follows from these measurements.

Saved failure records and three original images were inspected after evaluation.
Both models miss the annotated electronic region in \file{Test000001.jpg} at
confidence 0.25. The candidate produces a 0.945477 prediction on
\file{Test004918.jpg}, whose generic annotation export is empty. The known
modified-host label question in \file{Test000071.jpg} remains; both model scores
are below 0.25. These descriptive examples neither repair labels nor establish
benign outcomes, and no model or setting was changed after inspection.

The large validation-to-test AP gap must remain visible. Four original
train/test exact-byte overlaps and 112 within-test duplicate pairs were retained
to preserve the published split. Physical independence is unverified, one known
modified-host test image has a label question, and only one training seed was
used. No unseen-case, statistical significance or repeatability claim is made.
Specific classes and the two-stage host/threat or distributed pipelines were
not evaluated by this generic detector experiment.

\subsection{Timing and resource measurements}
Five validation-image warmups preceded each test pass. Across 5,136 timed images
per model, synchronized predictor-call median/95th-percentile times were
7.395/8.086 ms for the control and 7.387/7.886 ms for the candidate. This boundary
includes image loading, preprocessing, model execution and postprocessing;
input hashing and output writing are excluded. The fork's separate model-only
median/95th-percentile times were 5.589/5.713 ms and 5.583/5.691 ms respectively.
These measurements describe the allocated H100 NVL with batch one and FP32,
not GUI, scanner, network or complete operator-workflow latency. Slurm recorded
1,728,716 KiB maximum batch resident memory for job 820. Raw GPU samples,
checkpoint identities, prediction hashes and timings are retained privately.
\par\projectrecord{docs/t41-results.md} See also \file{docs/t41-handoff.md}.
